% GENERATED FILE. Edit canonical lessons, then run npm run generate:books.
% canonical-source-hash: fnv1a64:86fea7352f9bc080
% canonical-lessons: FR-C483-nausee, FR-C483-vomir, FR-C483-migraine, FR-C483-vertige, FR-C483-digerer

\chapter{Mal au cœur}
\label{ch:fr-mal-au-c-ur}
\hlchaptermodality{483}

\begin{chapteropening}
\textbf{By the end of this chapter:} \emph{I can describe feeling sick, dizzy or having a headache.}

Complete the last lesson of chapter 483: I can describe feeling sick, dizzy or having a headache.
\end{chapteropening}

\section[la nausée]{la nausée — nausea}

\label{lesson:FR-C483-nausee}



\begin{quote}
Before the new one: say the French for itch, then the French for to scratch.
\end{quote}

\subsection*{You'll want to know: la nausée}
\textbf{la nausée} — \textquotedblleft{}nausea\textquotedblright{}.

\begin{quote}
En bateau, j'ai souvent la nausée.
\end{quote}

\begin{quote}
\emph{On a boat I often feel sick.}
\end{quote}

\emph{Avoir la nausée} or, more informally, \emph{avoir mal au cœur}.

\begin{cousinweb}
\emph{la nausée} is a learned borrowing from Latin \emph{nausea} \textquotedblleft{}seasickness\textquotedblright{} (EtymDB), from Greek \emph{nausía}, built on \emph{naûs} \textquotedblleft{}ship\textquotedblright{}.
\end{cousinweb}

\begin{grammarlens}[title={What you've built}]
One more word for telling, explaining and describing at B1.
\end{grammarlens}

\subsection*{Your turn}
\begin{itemize}
\raggedright
  \item \emph{Say it:} \emph{la nausée}
  \item \emph{Say it:} \emph{la nausée}, once more
  \item \emph{Say it:} a sentence of your own with \emph{la nausée}
  \item \emph{Recall:} say the French for itch, then the French for to scratch, then say \emph{la nausée} again
\end{itemize}

\begin{tcolorbox}[breakable,colback=teal!4,colframe=teal!35!black,title={Before you move on}]
What does la nausée mean? (\textquotedblleft{}Nausea\textquotedblright{}.) Say it once more.
\end{tcolorbox}
\section[vomir]{vomir — to vomit, to be sick}

\label{lesson:FR-C483-vomir}



\begin{quote}
Before the new one: say the French for to scratch, then the French for nausea.
\end{quote}

\subsection*{You'll want to know: vomir}
\textbf{vomir} — \textquotedblleft{}to vomit, to be sick\textquotedblright{}.

\begin{quote}
Le petit a vomi dans la voiture.
\end{quote}

\begin{quote}
\emph{The little one was sick in the car.}
\end{quote}

A regular \emph{-ir} verb: \emph{nous vomissons}.

\begin{cousinweb}
\emph{vomir} is inherited from Vulgar Latin \emph{vomīre}, a form of Latin \emph{vomere} \textquotedblleft{}to vomit\textquotedblright{} (EtymDB; Lewis \& Short).
\end{cousinweb}

\begin{grammarlens}[title={What you've built}]
One more word for telling, explaining and describing at B1.
\end{grammarlens}

\subsection*{Your turn}
\begin{itemize}
\raggedright
  \item \emph{Say it:} \emph{vomir}
  \item \emph{Say it:} \emph{vomir}, once more
  \item \emph{Say it:} a sentence of your own with \emph{vomir}
  \item \emph{Recall:} say the French for to scratch, then the French for nausea, then say \emph{vomir} again
\end{itemize}

\begin{tcolorbox}[breakable,colback=teal!4,colframe=teal!35!black,title={Before you move on}]
What does vomir mean? (\textquotedblleft{}To vomit, to be sick\textquotedblright{}.) Say it once more.
\end{tcolorbox}
\section[la migraine]{la migraine — migraine, bad headache}

\label{lesson:FR-C483-migraine}



\begin{quote}
Before the new one: say the French for nausea, then the French for to vomit.
\end{quote}

\subsection*{You'll want to know: la migraine}
\textbf{la migraine} — \textquotedblleft{}migraine, bad headache\textquotedblright{}.

\begin{quote}
Quand j'ai la migraine, je reste dans le noir.
\end{quote}

\begin{quote}
\emph{When I have a migraine, I stay in the dark.}
\end{quote}

In speech, also any bad headache.

\begin{cousinweb}
\emph{la migraine} comes from Late Latin \emph{hēmicrānia} (EtymDB), from Greek \emph{hēmikranía}, built on \emph{hēmi-} \textquotedblleft{}half\textquotedblright{} and \emph{kraníon} \textquotedblleft{}skull\textquotedblright{}: a pain on one side of the head.
\end{cousinweb}

\begin{grammarlens}[title={What you've built}]
One more word for telling, explaining and describing at B1.
\end{grammarlens}

\subsection*{Your turn}
\begin{itemize}
\raggedright
  \item \emph{Say it:} \emph{la migraine}
  \item \emph{Say it:} \emph{la migraine}, once more
  \item \emph{Say it:} a sentence of your own with \emph{la migraine}
  \item \emph{Recall:} say the French for nausea, then the French for to vomit, then say \emph{la migraine} again
\end{itemize}

\begin{tcolorbox}[breakable,colback=teal!4,colframe=teal!35!black,title={Before you move on}]
What does la migraine mean? (\textquotedblleft{}Migraine, bad headache\textquotedblright{}.) Say it once more.
\end{tcolorbox}
\section[le vertige]{le vertige — dizziness; fear of heights}

\label{lesson:FR-C483-vertige}



\begin{quote}
Before the new one: say the French for to vomit, then the French for migraine.
\end{quote}

\subsection*{You'll want to know: le vertige}
\textbf{le vertige} — \textquotedblleft{}dizziness; fear of heights\textquotedblright{}.

\begin{quote}
En haut de la tour, j'ai eu le vertige.
\end{quote}

\begin{quote}
\emph{At the top of the tower, I felt dizzy.}
\end{quote}

\emph{Avoir le vertige} — to feel dizzy, or to be afraid of heights.

\begin{cousinweb}
\emph{le vertige} is a learned borrowing from Latin \emph{vertīgō} \textquotedblleft{}a turning round, dizziness\textquotedblright{}, from \emph{vertere} \textquotedblleft{}to turn\textquotedblright{} (Lewis \& Short).
\end{cousinweb}

\begin{grammarlens}[title={What you've built}]
One more word for telling, explaining and describing at B1.
\end{grammarlens}

\subsection*{Your turn}
\begin{itemize}
\raggedright
  \item \emph{Say it:} \emph{le vertige}
  \item \emph{Say it:} \emph{le vertige}, once more
  \item \emph{Say it:} a sentence of your own with \emph{le vertige}
  \item \emph{Recall:} say the French for to vomit, then the French for migraine, then say \emph{le vertige} again
\end{itemize}

\begin{tcolorbox}[breakable,colback=teal!4,colframe=teal!35!black,title={Before you move on}]
What does le vertige mean? (\textquotedblleft{}Dizziness; fear of heights\textquotedblright{}.) Say it once more.
\end{tcolorbox}
\section[digérer]{digérer — to digest}

\label{lesson:FR-C483-digerer}



\begin{quote}
Before the new one: say the French for migraine, then the French for dizziness.
\end{quote}

\subsection*{You'll want to know: digérer}
\textbf{digérer} — \textquotedblleft{}to digest\textquotedblright{}.

\begin{quote}
Je ne digère pas bien le lait.
\end{quote}

\begin{quote}
\emph{I do not digest milk well.}
\end{quote}

Also to take in bad news: \emph{il n'a pas digéré la défaite}.

\begin{cousinweb}
\emph{digérer} is a learned borrowing from Latin \emph{dīgerere} \textquotedblleft{}to divide, to arrange, to digest\textquotedblright{}, built on \emph{gerere} \textquotedblleft{}to carry\textquotedblright{} (Lewis \& Short).
\end{cousinweb}

\begin{grammarlens}[title={What you've built}]
One more word for telling, explaining and describing at B1.
\end{grammarlens}

\subsection*{Your turn}
\begin{itemize}
\raggedright
  \item \emph{Say it:} \emph{digérer}
  \item \emph{Say it:} \emph{digérer}, once more
  \item \emph{Say it:} a sentence of your own with \emph{digérer}
  \item \emph{Recall:} say the French for migraine, then the French for dizziness, then say \emph{digérer} again
\end{itemize}

\begin{tcolorbox}[breakable,colback=teal!4,colframe=teal!35!black,title={Before you move on}]
What does digérer mean? (\textquotedblleft{}To digest\textquotedblright{}.) Say it once more.
\end{tcolorbox}
